% ============================================
% 2 问题分析
% ============================================
\section{问题分析}

四个问题围绕"古代玻璃成分分析与鉴别"这一核心任务，在逻辑上层层递进，构成了从现象描述到规律挖掘的完整分析链条。

\subsection{问题1：风化规律与预测}

\textbf{问题性质}：统计推断与预测问题。

\textbf{分析要点}：风化是影响玻璃成分判断的主要干扰因素。该问题要求从三个层面逐步深入——首先定性判断风化与哪些因素（类型、纹饰、颜色）存在关联，然后定量分析风化前后各化学成分的变化幅度及统计显著性，最后建立数学模型将风化样品的成分"还原"至风化前状态。核心难点在于：（1）风化效应因玻璃类型不同而呈现截然不同的模式（高钾淋溶 vs 铅钡交换）；（2）校正向量的估计受限于配对样本数量。解决思路：用列联表$+$假设检验处理定性关联，用非参数检验$+$Mann-Whitney U 处理定量比较，用分类型差值法建模预测。

\textbf{方法选择依据}：1a 小问的"风化$\times$类型"为 $2 \times 2$ 列联表，期望频数偏小，选用 Fisher 精确检验而非卡方检验，确保小样本下 $p$ 值可靠；1b 小问的化学成分数据不服从正态分布（多数组分呈偏态、存在零值堆积），因此弃用 $t$ 检验和 ANOVA，改用 Mann-Whitney U 非参数检验配合 Bonferroni 多重校正，避免假阳性；1c 小问面临铅钡玻璃仅 2 对配对样本的约束，无法拟合多元回归或机器学习模型，故选用直观的差值校正法——其优势在于每个 $\Delta$ 值有明确的化学含义（流失/富集的方向和幅度），且留出验证证明关键氧化物预测误差 $< 1.5\%$，说明用简单模型解决小样本预测问题是合理且有效的。

\subsection{问题2：分类规律与亚类划分}

\textbf{问题性质}：分类与聚类问题。

\textbf{分析要点}：该问题需要在已知类别标签（高钾/铅钡）的基础上，一方面提炼出最简洁有效的分类规则，另一方面探索每类玻璃内部是否存在有意义的亚类结构。前者属于有监督学习的特征选择问题——找出能完美区分两类的化学指标；后者属于无监督学习的聚类问题——在未风化样品中识别配方多样性。核心难点在于：（1）分类规则必须对风化鲁棒，否则无法应用于风化样品；（2）亚类划分需区分"真实配方差异"和"风化程度差异"。解决思路：用 AUC/Cohen's $d$ 量化单氧化物判别力，用决策树提炼可解释规则，用 $\text{PCA} + K\text{-means}$ 做亚类聚类，用 Bootstrap 验证稳定性。

\textbf{方法选择依据}：分类器选决策树而非 LDA 或 SVM 的原因在于——LDA 在 25 个训练样本上依赖 SrO、$\text{SnO}_2$、MgO 等微量次要特征，这些特征在风化过程中不稳定，导致风化测试集准确率仅 $61.9\%$；而决策树仅选 BaO 单一特征（阈值 $3.14\%$），BaO 作为重元素助熔剂成分在风化中几乎不流失，风化测试集准确率仍保持 $90.5\%$。单特征规则另一个优势是化学可解释性强——BaO 是有意添加的助熔剂，其含量直接关联玻璃工艺类型，不是统计巧合。亚类划分选 $K\text{-means}$ 而非层次聚类或 DBSCAN 的原因在于——$K\text{-means}$ 配合肘部法则和轮廓系数可客观确定 $k$ 值，PCA 预降维（保留 $\geq 85\%$ 方差）消除了成分数据共线性对距离度量的干扰，Bootstrap 重采样提供了聚类稳定性的定量评估（ARI），使亚类结果可被检验而非主观判定。

\subsection{问题3：未知类别鉴别}

\textbf{问题性质}：分类应用与模型评估问题。

\textbf{分析要点}：该问题是将问题2建立的分类规则应用于 8 个未知样品（A1--A8），检验模型的泛化能力。由于部分未知样品为风化状态，分类器必须在风化干扰下保持准确。核心难点在于：（1）单一分类器可能存在偏差，需要多方法交叉验证；（2）未知样品中可能存在训练集未覆盖的新类型（如 A2），需要识别异常。解决思路：采用多分类器投票共识策略，对分类结果分置信度，对异常样品单独讨论。

\textbf{方法选择依据}：选用四法投票而非单一最优分类器的原因——虽然决策树在训练集上表现最好，但任何一个分类器都有盲区：LDA 对 A2（PbO 极高但 $\text{BaO}=0$）的分类与其它三法不一致，恰是这种分歧揭示了 A2 的异常性。若只用决策树单方法，A2 会被无争议地归为高钾，论文将错失"可能存在第三种玻璃类型"这一重要发现。不选用更复杂的集成模型（如 Stacking、随机森林）的原因——训练集仅 25 个样本且特征维度为 14，复杂模型容易过拟合；四个简单方法的投票结果已高度一致（BaO 规则与决策树 8/8 完全一致），说明简单方法足够。风化校正交叉验证的设计意图在于——即使分类规则对训练集的风化样品准确率 $90.5\%$，仍需要对未知风化样品做"校正后重分类"以排除残余风险，这是从问题1 到问题3 的严谨性闭环。

\subsection{问题4：化学成分关联关系}

\textbf{问题性质}：关联挖掘与比较分析问题。

\textbf{分析要点}：该问题要求从化学成分的共变关系出发，揭示玻璃制造工艺和风化过程的内在化学逻辑。由于成分数据存在闭合效应（总和固定为 $100\%$，一组分升高必然导致其他组分占比下降），直接计算 Pearson 相关系数会产生大量虚假负相关。核心难点在于：（1）必须在消除闭合效应的前提下计算真实关联；（2）需区分"原始配方关联"与"风化诱导关联"。解决思路：用 CLR 变换将成分数据从单形空间映射到欧氏空间，分别计算 Pearson 和 Spearman 相关矩阵，用 Fisher $z$ 检验比较类型间差异，按风化状态分层验证。

\textbf{方法选择依据}：选用 CLR 变换而非 ILR（等距对数比）变换的原因——CLR 变换后每个坐标仍对应一种原始氧化物，保留了化学可解释性（可以直接说"$\text{SiO}_2$ 与 $\text{SO}_2$ 在 CLR 空间中 $r = +0.80$"），而 ILR 坐标是原始成分的线性组合（如 $\ln(\text{SiO}_2/\text{某参考组分})$），结果无法直接对应到具体氧化物，论文可读性差。选用 Fisher $z$ 检验做类型间相关系数比较而非仅凭肉眼对比热力图的原因——视觉比较无法量化差异的统计显著性，Fisher $z$ 检验给出了严格的 $|z| > 2$ 判定标准，最终识别出 18/91 对氧化物的关联在两类玻璃间显著不同，这一数字本身就是一个可引用的核心发现。风化分层验证的设计意图——通过比较四组（高钾风化/未风化 $\times$ 铅钡风化/未风化）的平均 $|r|$，发现高钾风化组的平均 $|r|$ 异常高（$0.567$ vs 其他三组 $0.25\sim0.28$），直接证明了风化淋溶通过闭合效应制造了虚假高相关，从而验证了"风化分析必须分层"这一前提假设的正确性。

\subsection{问题间逻辑关系}

四个问题构成了完整的分析闭环：

\begin{center}
\begin{tabular}{cccc}
    问题1（风化分析） & $\rightarrow$ & 问题2（分类规律） & $\rightarrow$ \\
    $\downarrow$ & & $\downarrow$ & \\
    去除干扰因素 & & 建立判别规则 & \\
\end{tabular}
\begin{tabular}{cccc}
    问题3（未知鉴别） & $\rightarrow$ & 问题4（关联关系） \\
    $\downarrow$ & & $\downarrow$ \\
    检验模型泛化 & & 挖掘化学本质 \\
\end{tabular}
\end{center}

问题1 清除风化的"噪声"，为后续分析提供干净的数据基础；问题2 在问题1 的理解之上提炼出 BaO 等稳定判别指标；问题3 验证问题2 的分类规则在实际未知样品中的表现；问题4 回到化学本身，在问题1 的风化认知和问题2 的类型认知之上，揭示两类玻璃化学体系的本质差异。

\vspace{0.5cm}
\noindent\textbf{AI 工具使用说明}：本文在问题分析方法遴选与论证逻辑梳理过程中，借助了 DeepSeek-V4 大语言模型辅助分析，所有方法选择依据和推理结论均经人工核验确认。
